\documentclass[12pt,a4paper]{article}
\usepackage[utf8]{vietnam}
\usepackage{lmodern} % Font New Computer Modern chuẩn, thanh lịch
\usepackage{amsmath,amssymb}
\usepackage[margin=1.5cm]{geometry}
\usepackage{lastpage}

% ======================================================================
% CẤU HÌNH GÓI EX_TEST (BẬT/TẮT LỜI GIẢI DỄ DÀNG)
% ======================================================================
% [dethi]    : Dành cho học sinh (Chỉ in đề, ẩn đáp án)
% [loigiai]  : Dành cho giáo viên (In đề & lời giải chữ đen)
% [solcolor] : Dành cho giáo viên (In đề & lời giải chữ xanh)
% ======================================================================
\usepackage[dethi]{ex_test}  

% Gọi giao diện cao cấp Sang Math (Ẩn toàn bộ các lệnh cấu hình phức tạp)
\usepackage{sangmath-theme}

\begin{document}

% ======================================================================
% HEADER ĐỀ THI TỐI GIẢN (KHÔNG DƯ THỪA - RẤT CHUYÊN NGHIỆP)
% ======================================================================
\begin{center}
    \textbf{\textcolor{AccentColor}{\small TRƯỜNG THPT SANG MATH}}\\[5pt]
    \textbf{\Large ĐỀ KIỂM TRA TOÁN 12}\\[3pt]
    \textit{\small Thời gian làm bài: 45 phút \quad $\bullet$ \quad Trang 1/\pageref{LastPage}}
\end{center}
\vspace{2pt}
\noindent\textcolor{AccentColor}{\rule{\textwidth}{1pt}}
\vspace{10pt}

% ======================================================================
% NỘI DUNG CÂU HỎI
% ======================================================================
\begin{ex}
Tập xác định của hàm số $y = \frac{1}{x-1}$ là
\choice{$\mathbb{R}$}{\True $\mathbb{R} \setminus \{1\}$}{$(1; +\infty)$}{$\mathbb{R} \setminus \{0\}$}
\loigiai{
    Hàm số xác định khi $x - 1 \neq 0 \Leftrightarrow x \neq 1$. \\
    Vậy tập xác định là $D = \mathbb{R} \setminus \{1\}$.
}
\end{ex}

\begin{ex}
Nghiệm của phương trình $2^x=8$ là
\choice{$x = 1$}{$x = 2$}{\True $x = 3$}{$x = 4$}
\loigiai{
    Ta có $2^x = 8 \Leftrightarrow 2^x = 2^3 \Leftrightarrow x = 3$.
}
\end{ex}

\begin{ex}
Giá trị lớn nhất của hàm số $y=-x^2+4x$ bằng
\choice{$2$}{\True $4$}{$0$}{$-4$}
\loigiai{
    Hàm số bậc hai có hệ số $a=-1 < 0$ nên đạt giá trị lớn nhất tại đỉnh parabol. \\
    Hoành độ đỉnh $x = -\frac{b}{2a} = 2$. \\
    Suy ra $\max y = y(2) = -(2)^2 + 4(2) = 4$.
}
\end{ex}

\begin{ex}
Tính tích phân $\int_0^1 2x \, \mathrm{d}x$.
\choice{$0$}{\True $1$}{$2$}{$-1$}
\loigiai{
    Áp dụng công thức cơ bản:
    \[ \int_0^1 2x \, \mathrm{d}x = x^2 \Big|_0^1 = 1^2 - 0^2 = 1. \]
}
\end{ex}

\end{document}
